\BourbakiExerciseGroup

\begin{enumerate}
\def\labelenumi{\arabic{enumi})}
\item
  设 X 为拓扑空间，S 为 X 的拓扑的一个子基（§ 2，第 3 目）。证明：如果由属于 S 的集组成的 X 的每个开覆盖都含有 X 的一个有限覆盖，则 X 是拟紧的。（注意：若 X 上的超滤子 U 不收敛，则每个 \(x \in X\) 都属于一个不属于 U 的 S 中的集；并利用 § 6 命题 5 的推论。）
\item
  证明：良序集 X 赋以拓扑 \(\mathcal{C}_{-}(\mathbf{X})\)（§ 2，习题 5）后是局部紧的，且它是紧的当且仅当 X 有最大元。由此推出，每个集都可以赋以拓扑，使之成为一个紧空间。
\item
  \begin{enumerate}
  \def\labelenumii{\alph{enumii})}
  \tightlist
  \item
    设 X 为豪斯多夫空间，A、B 为 X 的两个不相交的紧子集。证明：存在 A 的一个邻域和 B 的一个邻域，二者没有公共点。
  \end{enumerate}
\end{enumerate}

\begin{enumerate}
\def\labelenumi{\alph{enumi})}
\setcounter{enumi}{1}
\tightlist
\item
  设 X 为正则空间，A 为 X 的紧子集，B 为 X 的不与 A 相交的闭子集。证明：存在 A 的一个邻域和 B 的一个邻域，二者没有公共点。
\end{enumerate}

\begin{enumerate}
\def\labelenumi{\arabic{enumi})}
\setcounter{enumi}{3}
\item
  证明：在 § 7 习题 6 所定义的空间 X（它是正则的（§ 8，习题 16））中，每个紧子集都是有限的。
\item
  *a) 证明：§ 8 习题 7 所定义的非豪斯多夫可达空间 Y = X/S 中，Y 的每个点都有紧邻域；1 与 -1 在 Y 中的像 α、β 有非闭的紧邻域；α 的一个紧邻域与 β 的一个紧邻域的交不是拟紧的，而这两个邻域的并不是紧的。*
\end{enumerate}

\begin{enumerate}
\def\labelenumi{\alph{enumi})}
\setcounter{enumi}{1}
\item
  设 X 是 N 与一个无限集 A 的和。对每个 \(x \in X\)，设 \(\mathfrak{S}(x)\) 是如下定义的滤子基：若 \(x \in N\)，则 \(\mathfrak{S}(x)\) 由单个集 \(\{x\}\) 组成；若 \(x \in A\)，且用 \(S_{x,n}\) 表示以 x 及整数 \(\geqslant n\) 为元素的集，则 \(\mathfrak{S}(x)\) 由诸集 \(S_{x,n}\)（n 取遍 N）组成。证明：X 上存在一个拓扑，使得 \(\mathfrak{S}(x)\) 对每个 \(x \in X\) 都是 x 的基本邻域系，且在此拓扑下 X 是可达的但不是豪斯多夫的。再证明 X 不是拟紧的，但 X 有稠密的紧子集。
\item
  证明：集 X 上的两个拓扑，若 X 在其每个之下都是可达的且拟紧的，则它们可以是不同但可比较的{[}参见§ 8，习题 7 及 5 b){]}。
\end{enumerate}

\begin{enumerate}
\def\labelenumi{\arabic{enumi})}
\setcounter{enumi}{5}
\item
  设 X、Y 为两个拓扑空间，A（相应地，B）为 X（相应地，Y）的拟紧子集。证明：若 U 是 \(A \times B\) 在 \(X \times Y\) 中的任一邻域，则存在 X 中 A 的邻域 V 与 Y 中 B 的邻域 W，使得 \(V \times W \subset U\)。
\item
  \begin{enumerate}
  \def\labelenumii{\alph{enumii})}
  \tightlist
  \item
    试给出拟紧空间的一个逆系统 \((\mathbf{X}_{\alpha}, f_{\alpha\beta})\)，它的逆极限是空集（注意：在第 1 目例 2 的空间中，每个子空间都是拟紧的）。
  \end{enumerate}
\end{enumerate}

\begin{enumerate}
\def\labelenumi{\alph{enumi})}
\setcounter{enumi}{1}
\tightlist
\item
  设 \(p\) 为素数。用 \(Z_{n}\) 表示有理整数集 \(\mathbf{Z}\)，赋以 \(\mathbf{Z}\) 到 \(\mathbf{Z} / p^{n}\mathbf{Z}\) 上的典范映射下离散拓扑的逆像。若 \(m \leqslant n\)，用 \(f_{mn}\) 表示恒等映射 \(Z_{n} \to Z_{m}\)，它是连续的。证明：诸空间 \(Z_{n}\) 都是拟紧的，但逆系统 \((Z_{n}, f_{mn})\) 的逆极限不是拟紧的（*）。
\end{enumerate}

\begin{enumerate}
\def\labelenumi{\arabic{enumi})}
\setcounter{enumi}{7}
\tightlist
\item
  设 X 为拓扑空间，\((\mathrm{G}_{\alpha})_{\alpha\in\mathbb{A}}\) 为 X 上等价关系图像的一个族，其指标集 A 是有向的；设 \(R_{\alpha}\) 为关系 \((x,y)\in\mathbf{G}_{\alpha}\)，\(\varphi_{\alpha}\) 为典范映射 \(X\to X/R_{\alpha}\)。设关系 \(R_{\beta}\) 蕴含 \(R_{\alpha}\)（只要 \(\alpha\leqslant\beta\)），并设 \(\varphi_{\alpha\beta}\) 为典范映射 \(X/R_{\beta}\to X/R_{\alpha}\)；于是 \((X/R_{\alpha},\varphi_{\alpha\beta})\) 是拓扑空间的一个逆系统。设 \(Y=\varprojlim X/R_{\alpha}\)。
\end{enumerate}

\begin{enumerate}
\def\labelenumi{\alph{enumi})}
\item
  设 \(g: \mathbf{X} \to \mathbf{Y}\) 为连续映射 \(\varprojlim \varphi_{\alpha}\)。证明：\(g(x)\) 在 \(\mathbf{Y}\) 中稠密，且对每个 \(x \in \mathbf{X}\)，\(\overrightarrow{g}(g(x))\) 是 \(x\) 关于其图像为 \(\bigcap_{n} G_{\alpha}\) 的等价关系的类。
\item
  设所有 \(R_{\alpha}\) 都是闭等价关系，且对每个 \(x \in X\) 与 x 的每个邻域 V，存在指标 \(\alpha\)，使得 x 关于 \(R_{\alpha}\) 的类含于 V。证明：在这些条件下，g 是 X 到 \(g(\mathbf{X})\) 上的开映射。
\item
  证明：若模 \(R_{\alpha}\) 的各类对每个 \(\alpha \in A\) 都是闭的且拟紧的，则 g 是满射。
\end{enumerate}

\begin{enumerate}
\def\labelenumi{\arabic{enumi})}
\setcounter{enumi}{8}
\item
  设 X 为可达空间，\(\mathcal{R}\) 为 X 的拟紧闭子集所成的一个集，它包含 X 的所有有限子集，且 \(\mathcal{R}\) 的任意两个子集之并属于 \(\mathcal{R}\)。设 \(X'\) 为一个集，它是 X 与由单点组成的集 \(\{\omega\}\) 的和。证明：\(X'\) 上存在一个拟紧的可达拓扑，它在 X 上诱导出给定的拓扑，且使得 \(X'\) 中 \(\mathcal{R}\) 的诸集之余集构成 \(\omega\) 的开邻域的一个基本系。试给出 \(X'\) 上与不同的集 \(\mathcal{R}\) 相对应的不同拓扑的例子。
\item
  设 X 为仿紧空间。证明：若 X 的每个开子空间都是仿紧的，则 X 的每个子空间都是仿紧的。由此推出：拓扑具有可数基的局部紧空间，其每个子空间都是仿紧的{[}参见第九章，§ 4，习题 20 a){]}。
\end{enumerate}

¶ 11) 设 \(\mathbf{X}_0\) 为具有最大元的不可数良序集；设 \(a\) 为 \(\mathbf{X}_0\) 的最小元，\(b\) 为元素 \(x \in \mathbf{X}_0\) 中使区间 \([a, x[\) 不可数的最小者。用 \(\mathbf{X}\) 表示区间 \([a, b[\)，赋以拓扑 \(\mathfrak{C}_{-}(\mathbf{X})\)。

\begin{enumerate}
\def\labelenumi{\alph{enumi})}
\item
  X 是局部紧的但不是紧的（习题 2），且 X 中每个序列都有聚点（注意 X 的每个可数子集都是上有界的）。
\item
  设 \(f\) 是 \(\mathbf{X}\) 到自身的映射，使得 \(f(x) < x\) 对一切充分大的 \(x\) 成立。证明：存在元素 \(c \in \mathbf{X}\)，具有如下性质：对每个 \(x \in \mathbf{X}\)，存在 \(y \geqslant x\) 使得 \(f(y) \leqslant c\)。{[}用反证法，构造点序列 \((z_n)\)（\(\mathbf{X}\) 的点），使得 \(z_{n+1}\) 是元素 \(z' \in \mathbf{X}\) 中满足 \(f(x) \geqslant z_n\)（对每个 \(x \geqslant z'\)）的最小者{]}。
\item
  由 b) 推出 X 不是仿紧的。
\end{enumerate}

\begin{enumerate}
\def\labelenumi{\arabic{enumi})}
\setcounter{enumi}{11}
\tightlist
\item
  设 X 为拓扑空间，\(\mathfrak{F}(X)\) 为 X 的非空闭子集所成之集，\(\mathfrak{K}(X)\) 为 X 的非空拟紧子集所成之集；
\end{enumerate}

给这些集中的每一个赋以 \(\mathfrak{C}_{\Omega}\)（\(\mathfrak{P}_{0}(\mathbf{X})\) 上的拓扑）所诱导的拓扑（§ 2，习题 7）。

\begin{enumerate}
\def\labelenumi{\alph{enumi})}
\tightlist
\item
  证明：若 \(\mathfrak{B} \subset \mathfrak{F}(\mathbf{X})\) 关于 \(\mathcal{C}_{\Omega}\) 是拟紧的，且 \(\mathbf{X}\) 是正则的，则诸集 \(\mathbf{M} \in \mathfrak{B}\) 的并在 \(\mathbf{X}\) 中是闭的。\\
\item
  证明：若 \(\mathfrak{B} \subset \mathfrak{K}(\mathbf{X})\) 关于 \(\mathcal{C}_{\Omega}\) 是拟紧的，则诸集 \(\mathbf{M} \in \mathfrak{B}\) 的并是拟紧的。
\end{enumerate}

¶ 13) 用习题 13 的记号，给 \(\mathfrak{F}(\mathbf{X})\) 和 \(\mathfrak{R}(\mathbf{X})\) 赋以 \(\mathfrak{C}_{\Theta}\)（\(\mathfrak{P}_{0}(\mathbf{X})\) 上的拓扑）所诱导的拓扑（§ 8，习题 12）。a) 证明：若 X 是拟紧的，则 \(\mathfrak{F}(\mathbf{X})\) 也是。（对 X 的每个开集 U，设 \(c(\mathrm{U})\){[}相应地 \(m(\mathrm{U})\){]}为 X 中所有满足 \(\mathrm{M} \subset \mathrm{U}\)（相应地 \(\mathrm{M} \cap \mathrm{U} \neq \emptyset\)）的闭子集 M 所成之集；诸 \(c(\mathrm{U})\) 与 \(m(\mathrm{U})\) 构成 \(\mathfrak{C}_{\Theta}\) 的一个子基。然后利用习题 1）。

\begin{enumerate}
\def\labelenumi{\alph{enumi})}
\setcounter{enumi}{1}
\item
  设 X 是可达的，并设 \(X'\) 是 X 在映射 \(x \to \{x\}\)（X 到 \(\mathfrak{F}(X)\) 中的映射）下的像。设 H 为 \(\mathfrak{F}(X)\) 的一个包含 \(X'\) 的子空间。证明：若 H 是拟紧的，则 X 也是。{[}若 \((\mathrm{U}_{\alpha})\) 是 X 的一个开覆盖，则注意：若用 \(B_{\alpha}\) 表示所有 \(M \in \mathfrak{F}(X)\) 中满足 \(M \cap U_{\alpha} \neq \emptyset\) 者所成之集，则 \(\{\mathfrak{B}_{\alpha}\}\) 是 \(\mathfrak{F}(X)]\) 的一个开覆盖。
\item
  证明：若 X 是局部紧的，则 \(\mathfrak{K}(\mathbf{X})\) 在 \(\mathfrak{F}(\mathbf{X})\) 中是开的（利用 § 7 命题 10）。
\item
  设 X 是可达的。证明：\(\mathcal{R}(X)\) 是豪斯多夫的（相应地，正则的、紧的、局部紧的）当且仅当 X 如此。{[}对前两个断言，利用 § 8 习题 12 与 § 9 习题 3；对后两个断言，利用 a)、b) 及习题 12{]}。
\end{enumerate}

¶ 14) 拓扑空间 X 称为 Lindelöf 空间，如果 X 的每个开覆盖都含有 X 的一个可数覆盖。每个具有可数基的空间都是 Lindelöf 空间，每个拟紧空间也是。

\begin{enumerate}
\def\labelenumi{\alph{enumi})}
\item
  Lindelöf 空间的每个闭子空间都是 Lindelöf 空间。为使 Lindelöf 空间的每个子空间都是 Lindelöf 空间，必须且只需每个开子空间都是 Lindelöf 的（参见习题 15）。
\item
  若 \(f\) 是 Lindelöf 空间 \(\mathbf{X}\) 到拓扑空间 \(\mathbf{X}'\) 中的任一连续映射，则子空间 \(f(\mathbf{X})\)（\(\mathbf{X}'\) 的）是 Lindelöf 的。
\item
  拓扑空间的 Lindelöf 子空间的每个可数并都是 Lindelöf 的。特别地，每个可数空间都是 Lindelöf 的。
\item
  Lindelöf 空间 X 若 X 中每个序列都有聚点，则是拟紧的{[}证明此时公理 (C'') 成立{]}。
\end{enumerate}

¶ 15) a) 设 X 为每个开子空间都是 Lindelöf 的拓扑空间。证明：X 满足 § 1 习题 7 中的条件 (D \(_{III}\) )，且若 X 是正则的，则 X 的每个闭子集都是开集的可数交{[}参见习题 23 c){]}。

\begin{enumerate}
\def\labelenumi{\alph{enumi})}
\setcounter{enumi}{1}
\item
  设 X 为每个开子空间都是仿紧的拓扑空间。证明：X 的每个开子空间都是 Lindelöf 的，当且仅当 X 满足 §1 习题 7 的条件 \((D_{\mathrm{III}})\)。
\item
  设 X 为紧空间。证明：X 的每个开子空间都是 Lindelöf 的，当且仅当 (i) X 满足 §1 习题 7 的条件 \((\mathbf{D}_{\mathbf{III}})\)，且 (ii) X 的每个闭子集都是开集的可数交（利用 b) 与第 10 目定理 5）。
\end{enumerate}

¶ 16) 设 X 为拓扑空间，A 为 X 的子集。点 \(a \in X\) 称为 A 的凝聚点，如果 \(a\) 的每个邻域都含有 A 的不可数无穷多个点。

\begin{enumerate}
\def\labelenumi{\alph{enumi})}
\item
  X 的子集 A 的凝聚点所成之集是闭的。试给出一个例子，其中 X 只有一个凝聚点（参见第 8 目定理 4）。
\item
  设 X 是每个开子空间都是 Lindelöf 的可达空间（习题 15）。证明：对 X 的每个子集 A，A 的凝聚点所成之集 B 是完备集，且 A ∩ GB 是可数的。
\end{enumerate}

¶ 17) 在拓扑空间 X 中，点 \(a\) 称为与子集 A 全然附着，如果对 \(a\) 的每个邻域 U 都有 card (A) = card (A ∩ U)。证明：X 是拟紧的，当且仅当对 X 的每个无限子集 A，存在 X 的一点与 A 全然附着。（若 X 不是拟紧的，证明存在初始序数 \(\omega_{\alpha}\)（《集合论》第三章，§ 6，习题 10）以及 X 的开覆盖 (U \(_{\xi}\) )，其中 \(\xi\) 取遍序数集 \(<\omega_{\alpha}\)，使得对每个指标 \(\xi<\omega_{\alpha}\)，集 A \(_{\xi}\)（诸 U \(_{\eta}\)（\(\eta<\xi\)）之并在 X 中的余集）的势等于 N \(_{\alpha}\)；由此推出，可用超限归纳法定义 X 的点的一个族 ( \(x_{\xi}\) ) ( \(\xi<\omega_{\alpha}\) )，使得 \(x_{\xi}\in\mathrm{A}_{\xi}\) 对每个 \(\xi\) 成立，且 \(x_{\xi}\neq x_{\eta}\) 每当 \(<\eta\xi\)。）

¶ 18) 豪斯多夫空间 X 称为绝对闭的，如果它具有下述性质：对 X 到豪斯多夫空间 X' 的子空间上的每个同胚 f，f(X) 在 X' 中是闭的。每个紧空间都是绝对闭的。

\begin{enumerate}
\def\labelenumi{\alph{enumi})}
\tightlist
\item
  证明下列条件等价：
\end{enumerate}

(AF) 空间 X 是绝对闭的。

(AF') X 上由开集组成的每个滤子基至少有一个聚点。

(AF'') X 中交为空的闭子集族含有这样的有限子族：该子族中诸子集的内部之交为空。

(AF''') X 的每个开覆盖都含有诸闭包能覆盖 X 的有限子族。

\begin{enumerate}
\def\labelenumi{\alph{enumi})}
\setcounter{enumi}{1}
\item
  设 X 为豪斯多夫空间，U 为 X 的稠密开子集。设 U 上由开集组成的每个滤子基在 X 中至少有一个聚点。证明：在这些条件下 X 是绝对闭的。特别地，若 A 是绝对闭空间 X 的开子集，则 \(\overline{A}\) 是绝对闭的。
\item
  设 X 为绝对闭空间，f 为 X 到豪斯多夫空间 \(X'\) 中的连续映射。证明：\(f(\mathbf{X})\) 是 \(X'\) 的绝对闭子空间。
\item
  证明：两个绝对闭空间的积是绝对闭的（按第 10 目命题 17 的方式论证）。
\end{enumerate}

¶ 19) 豪斯多夫空间 X 称为极小的，如果它的拓扑 C 满足：X 上每个豪斯多夫且比 C 更粗的拓扑必与 C 相合；此时 C 必是半正则的（§ 8，习题 20）。证明：豪斯多夫空间 X 是极小的，当且仅当 X 上每个由开集组成且至多有一个聚点的滤子基都是收敛的。这相当于说：X 是绝对闭的（习题 18），且 X 上每个由开集组成且有单个聚点的滤子基都收敛于该点。

¶ 20) 拓扑空间 X 称为完全豪斯多夫的（且 X 的拓扑称为完全豪斯多夫的），如果 X 的每对不同的点都有不相交的闭邻域。

\begin{enumerate}
\def\labelenumi{\alph{enumi})}
\item
  每个正则空间都是完全豪斯多夫的。完全豪斯多夫空间的每个子空间都是完全豪斯多夫的。完全豪斯多夫空间的每个积都是完全豪斯多夫的。比完全豪斯多夫拓扑更细的每个拓扑都是完全豪斯多夫的。
\item
  每个极小的完全豪斯多夫空间 X（习题 19）都是紧的。（用反证法：考虑 X 上的超滤子 U 以及由属于 U 的开集所组成的滤子基）。
\end{enumerate}

\(^{*}c)\) 设 \(\theta\) 为无理数 \textgreater0。对每个点 \((x,y)\in\mathbf{Q}^{2}\) 及每个 \(\alpha>0\)，设 \(\mathrm{B}_{\alpha}(x,y)\) 为由 \((x,y)\) 与点 \((z,0)\in\mathbf{Q}^{2}\)（满足 \(|z-(x+\theta y)|<\alpha\) 或 \(|z-(x-\theta y)|<\alpha\)）所组成的集。设 \(\mathfrak{B}(x,y)\) 为诸 \(\mathrm{B}_{\alpha}(x,y)\)（\(\alpha\) 取遍实数集 \textgreater0）所成之集。证明：存在豪斯多夫拓扑 \(\mathcal{C}\)（在 \(\mathbf{Q}^{2}\) 上），使得对每个 \((x,y)\)，\(\mathfrak{B}(x,y)\) 都是 \((x,y)\) 关于 \(\mathcal{C}\) 的基本邻域系。证明：对 \(\mathbf{Q}^{2}\) 的任意两点 a、b，在拓扑 \(\mathcal{C}.*\) 下，a 的每个闭邻域都与 b 的每个闭邻域相交。

¶ 21) a) 设 X 为豪斯多夫空间，\(\mathfrak{G}\) 为 X 的拓扑。证明：X 是绝对闭的（习题 18）当且仅当半正则拓扑 \(\mathfrak{G}^*\)（与 \(\mathfrak{G}\) 相伴随）（ \(\S\) 8，习题 20）是某个极小空间的拓扑（习题 19）。

\begin{enumerate}
\def\labelenumi{\alph{enumi})}
\setcounter{enumi}{1}
\tightlist
\item
  若 \(\mathfrak{C}\) 是完全豪斯多夫的（习题 20），则 \(\mathfrak{C}_{*}\) 也是（利用 § 8 习题 20）。由此推出，绝对闭的正则空间是紧的。
\end{enumerate}

¶ 22) a) 设 X 为豪斯多夫空间，A 为 X 的稠密子集。证明：若 A 上的每个滤子都在 X 中至少有一个聚点，则 X 是绝对闭的{[}验证习题 18 的公理 (AF'){]}。

\begin{enumerate}
\def\labelenumi{\alph{enumi})}
\setcounter{enumi}{1}
\tightlist
\item
  设 \(X_{0}\) 为具有非开稠密子集 A 的紧空间，\(C_{0}\) 为 \(X_{0}\) 的拓扑。设 C 为 \(X_{0}\) 上由 A 与 \(X_{0}\) 的在 \(C_{0}\) 中为开的子集所生成的拓扑，于是 \(C^{*} = C_{0}\)（§ 8，习题 20）。若用 X 表示把集 \(X_{0}\) 赋以拓扑 C（习题 21）所得的非紧绝对闭空间，证明 A 上的每个滤子在 X 中都有聚点。
\end{enumerate}

¶ 23) a) 设 X 为可数的绝对闭空间。证明：X 的孤立点所成之集在 X 中稠密。{[}用反证法：把 X 的点排成一个序列，并用归纳法构造 X 的可数开覆盖 (U \(_{n}\) )，使得诸 \(\overline{U}_{n}\) 的任何有限并不覆盖 X{]}。

*b) 试给出一个不是 Lindelöf 空间的完全豪斯多夫绝对闭拓扑空间的例子{[}在积 \([0,1]\times[0,1]\) 中，考虑诸集 \(A\times\{o\}\) 的余集，其中 A 取遍 \([0,1]\) 的所有子集，并应用 § 8 习题 20 c) 的一般方法{]}。这个空间满足 § 1 习题 1 的条件 \((D_{II})\)，但不满足条件 \((D_{III})\)。

\begin{enumerate}
\def\labelenumi{\alph{enumi})}
\setcounter{enumi}{2}
\tightlist
\item
  设 \(X_{0}\) 为没有孤立点的紧空间，M 为 \(X_{0}\) 的可数子集的余集所成之集{[}于是 M 中的集都在 \(X_{0}\) 中稠密，参见 a){]}。设 C 为由 \(X_{0}\) 的开集与 M 中的集所生成的拓扑。证明：把 \(X_{0}\) 赋以拓扑 C 所得的空间 X 是绝对闭的且完全豪斯多夫的，并且是一个不满足 §1 习题 7 条件 ( \(D_{II}\) ) 的 Lindelöf 空间；再证明 X 的任何点都没有可数的基本邻域系，且 X 中任何各项皆不同的序列都没有聚点。若 \(X_{0}\) 的每个开子空间都是 Lindelöf 的{[}参见习题 16 c){]}，证明 X 亦然，因而特别地 X 满足 ( \(D_{III}\) ){[}习题 16 a){]}。当 \(X_{0} = [0,1]\) 时即属此情形。但证明在此情形下，存在 X 的可数闭子集，它们不是 X 的开集的可数交（参见第九章，§5，第 3 目）。*
\end{enumerate}

¶ 24) a) 设 X 为豪斯多夫空间，Γ 为 X 上这样的滤子全体所成之集：它们具有由开集组成的基，且在 X 中没有聚点。证明：集 Γ 若非空（即若 X 不是绝对闭的），则是归纳的。

\begin{enumerate}
\def\labelenumi{\alph{enumi})}
\setcounter{enumi}{1}
\tightlist
\item
  设 \(\Omega\) 为 \(\Gamma\) 中极大滤子所成之集，并设 \(X'\) 为 \(X\) 与 \(\Omega\) 的和所成之集。对每个 \(x \in X\)，用 \(\mathfrak{B}(x)\) 表示 \(x\) 在 \(X\) 中的邻域所成之集；对每个 \(\omega \in \Omega\)，用 \(\mathfrak{B}(\omega)\) 表示 \(X'\) 的形如 \(\{\omega\} \cup M\) 的子集所成之集，其中 \(M\) 属于滤子 \(\omega\)。于是 \(X'\) 上存在一个拓扑，使得 \(\mathfrak{B}(x')\) 是 \(x'\)（对每个 \(x' \in X'\)）的基本邻域系。证明：在此拓扑下，\(X'\) 是绝对闭的，且具有下述泛性质：对任何开连续映射 \(f: X \to Y\)（\(X\) 到绝对闭空间 \(Y\) 中的映射），存在唯一的连续映射 \(g: X' \to Y\) 延拓 \(f\)（考察诸滤子在 \(f\) 下的像，这些滤子属于 \(\Omega\)，并注意在绝对闭空间中，在具有由开集组成的基的滤子集中为极大的滤子必收敛）。特别地，\(X'\) 的拓扑是 X-极大的（ \(\S 3\) ，习题 11）。由此推出，存在拓扑为拟极大的绝对闭空间（ \(\S 2\) ，习题 6）。
\end{enumerate}

¶ 25) a) 设 X 为可达空间，\(\Phi\) 为 X 上具有由闭集组成的基的滤子所成之集。证明：集 \(\Phi\) 是归纳的，并设 \(X''\) 为 \(\Phi\) 中极大滤子所成之集。

\begin{enumerate}
\def\labelenumi{\alph{enumi})}
\setcounter{enumi}{1}
\tightlist
\item
  对 X 的每个开子集 U，设 \(U^{*}\) 为滤子 \(F \in X''\)（\(U \in F\) 者）所成之集。若 U、V 为 X 的两个开子集，则
\end{enumerate}

\[
(\mathrm{U} \cap \mathrm{V}) ^ {*} = \mathrm{U} ^ {*} \cap \mathrm{V} ^ {*} \quad \text { and } \quad (\mathrm{U} \cup \mathrm{V}) ^ {*} = \mathrm{U} ^ {*} \cup \mathrm{V} ^ {*}
\]

（注意：若 \(\mathfrak{F} \in \mathbf{X}''\) 满足 \(\mathbf{U} \notin \mathfrak{F}\)，则 \(\mathbf{X} = \mathbf{U} \in \mathfrak{F}\)）。证明：诸集 \(\mathbf{U}^*\) 构成 \(\mathbf{X}''\) 上一个拓扑的基，且在此拓扑下 \(\mathbf{X}''\) 是拟紧的{[}利用上述注意，验证公理 \((\mathbf{C}''')\){]}并且是可达的{[}若 \(\mathfrak{F}, \mathfrak{G}\) 是 \(\mathbf{X}''\) 的两个不同元素，证明存在两个不相交的闭子集 \(\mathbf{A}\) 与 \(\mathbf{B}\)（\(\mathbf{X}\) 中的），使得 \(\mathbf{A} \in \mathfrak{F}\) 且 \(\mathbf{B} \in \mathfrak{G}\){]}。

\begin{enumerate}
\def\labelenumi{\alph{enumi})}
\setcounter{enumi}{2}
\item
  对每个 \(x \in \mathbf{X}\)，设 \(\varphi(x)\) 为 \(\mathbf{X}\) 上以单个集 \(\{x\}\) 为基的超滤子。证明：\(\varphi\) 是 \(\mathbf{X}\) 到 \(\mathbf{X}''\) 的一个稠密子集上的同胚，且若 \(\mathbf{X}\) 是拟紧的（且可达的），则 \(\varphi(\mathbf{X}) = \mathbf{X}''\)。
\item
  证明：对每个连续映射 \(f\)（\(\mathbf{X}\) 到紧空间 \(\mathbf{Y}\) 中的映射），存在唯一的连续映射 \(g: \mathbf{X}'' \to \mathbf{Y}\)，使得 \(f = g \circ \varphi\)。
\item
  设 \(\tilde{\mathbf{X}}\) 为与 \(\mathbf{X}''\) 相伴随的泛豪斯多夫空间（§ 8，习题 27）；证明：典范映射 \(\psi: \mathbf{X}'' \to \tilde{\mathbf{X}}\) 是满射，且 \(\tilde{\mathbf{X}}\) 是紧的。于是，\(\mathbf{X}\) 到紧空间 \(\mathbf{Y}\) 中的每个连续映射都可以唯一地写成 \(f = h \circ \psi \circ \varphi\) 的形式，其中 \(h: \tilde{\mathbf{X}} \to \mathbf{Y}\) 是连续的。
\end{enumerate}

\begin{enumerate}
\def\labelenumi{\arabic{enumi})}
\setcounter{enumi}{25}
\tightlist
\item
  若 X 是离散空间，证明习题 25 所定义的空间 \(X''\) 与 \(\tilde{X}\) 相同，且若 C 是习题 24 所定义的空间 \(X'\) 的拓扑，则 \(X''\) 可等同于空间 \(X'\)（赋有与 C 相伴随的半正则拓扑 \(C^*\)）（§ 8，习题 20）。此外，集 \(X''\) 可典范地等同于 X 上的超滤子所成之集。紧空间 \(X''\) 称为 X 的超滤子空间。
\end{enumerate}

¶ 27) 证明：在紧空间 X 中，每个等势的{[}§ 2，习题 11 d){]}无限开子空间都是可解的。{[}利用第 2 目命题 1，证明存在 U 的拓扑的一个基 B，使得 card (B) = card (U)；然后利用 § 2 习题 11 c)。{]}由此推出，每个没有孤立点的局部紧空间都是可解的{[}利用 § 2 习题 11 d){]}；因而这样的空间不是次极大的。

\begin{enumerate}
\def\labelenumi{\arabic{enumi})}
\setcounter{enumi}{27}
\tightlist
\item
  \begin{enumerate}
  \def\labelenumii{\alph{enumii})}
  \tightlist
  \item
    设 K 为离散空间 \(\{0,1\}\)。若 X 是任一紧空间，设 F 为 K 到 X 的闭子集所成之集 \(\mathfrak{F}(X)\) 中的映射 f 全体所成之集，这些映射满足 \(\mathbf{X}=f(0)\cup f(1)\)。考察紧空间 \(Y=K^{F}\)；证明：对 Y 中的每个 \(y=(y_{f})_{f\in\mathbb{F}}\)，X 的闭集 \(f(y_{f})\) 之交至多由一个点组成。使此交非空的 \(y\in Y\) 所成之集 Z 是 Y 的闭子集。若 \(\varphi(y)\) 是诸 \(f(y_{f})\)（对 \(y\in Z\)）公共的唯一点，证明 \(\varphi\) 是 Z 到 X 上的连续满射。
  \end{enumerate}
\end{enumerate}

\begin{enumerate}
\def\labelenumi{\alph{enumi})}
\setcounter{enumi}{1}
\tightlist
\item
  试给出这样的紧空间 X 的例子：不存在形如 \(K^{A}\) 的紧空间到 X 上的连续满射{[}利用 § 4 习题 4 b) 与 § 9 习题 26{]}。
\end{enumerate}

§ 29) 拓扑空间 X 称为局部拟紧的，如果 X 的每个点都有拟紧邻域的一个基本系。a) 证明：在局部拟紧空间中，每个拟紧子集都有拟紧邻域的一个基本系。

\begin{enumerate}
\def\labelenumi{\alph{enumi})}
\setcounter{enumi}{1}
\tightlist
\item
  证明：局部拟紧空间的每个局部闭子空间都是局部拟紧的。
\end{enumerate}

\(^{*}c)\) 在闭欧几里得圆盘 B：\(||x|| \leqslant 1\)（实平面 \(R^{2}\) 中的）上，设 \(\mathcal{G}_{0}\) 为由 \(R^{2}\) 的拓扑诱导的拓扑，并设 \(\mathcal{G}\) 为如下定义的（更细的）拓扑：B 的点 \(x \neq (1,0)\) 的邻域滤子在 \(\mathcal{G}\) 下与在 \(\mathcal{G}_{0}\) 下相同，但对点 \((1,0)\)，拓扑 \(\mathcal{G}\) 中的基本邻域系由诸集

\[
\left\{\left(\mathrm{I}, \mathrm{o}\right) \right\} \cup \left(\mathbf {D} \cap \mathbf {B} _ {\mathbf {0}}\right),
\]

组成，其中 D 是以 (1,0) 为中心的任一开圆盘，而 \(B_{0}\) 是开圆盘 \(\|x\|<1\)。设 Y 是赋以拓扑 C 的集 B，并设 R 是 Y 上的等价关系，其等价类是集 S——由所有 \(x\neq(1,0)\)

（B 中使 \(\|x\| = 1\) 的点）组成——以及由 B - S 的单个点组成的 Y 的子集。证明：商空间 X = Y/R 是拟紧的但不是局部拟紧的。*

